\documentclass{situationist-preprint}

\title{Intentionality Without Exception}
\msid{PHE-04-a}
\msfamily{PHE-04}
\msdiscipline{Phenomenology}
\mstype{Principal manuscript}
\mskeywords{intentionality; Brentano; Husserl; moods; profound boredom; Heidegger; time}

\begin{document}
\maketitle

\begin{abstract}
Is consciousness always consciousness of something? Brentano's thesis that intentionality is the mark of the mental faces counterexamples: sensations, undirected moods, and above all the profound boredom that Heidegger describes, a state defined by indifference to every being and every activity. We argue that the thesis holds without exception. A conscious state is intentional if it presents something as being some way; sensations present parts of the body, and undirected moods present the world or one's situation as a whole. Profound boredom is the hardest case, because what it presents is neither a being nor the world of beings, which withdraw in it as a whole. Following Heidegger's analysis of its three forms, we show that it nonetheless presents something as being some way: it presents time, the time of one's existence, as dragging and empty. Profound boredom is therefore intentional towards time itself, not towards an activity. The result distinguishes it from the ordinary boredom of tasks and from anxiety, which presents the nothing.
\end{abstract}

\section{The problem}

\citet[bk.~II, ch.~1]{Brentano1874} proposed that every mental phenomenon contains an object within itself: in presentation something is presented, in judgment something affirmed or denied, in love loved, in hate hated. Husserl took over the thesis in a corrected form, abandoning the language of immanent objects: consciousness is consciousness of something, in the sense that each act has a directedness determined by its matter and quality, whether or not its object exists \citep[V]{Husserl1901}. Intentionality became the theme of phenomenology. The thesis that it is universal, that there is no consciousness that is not consciousness of something, has been doubted on three grounds.

First, sensations such as pains, tickles and feelings of warmth seem to be states rather than presentations of objects; \citet{Searle1983} counts them as non-intentional. Second, moods differ from emotions in lacking a particular object: one is afraid of a dog, but one may be elated, depressed or anxious about nothing in particular \citep{Searle1983, Goldie2000}. Third, and most radically, Heidegger described a profound boredom in which every being and every activity becomes indifferent, so that the bored subject seems to be conscious and yet conscious of nothing \citep[\S\S 18--38]{Heidegger1995}. If any state is a counterexample, this is it, since it is defined by the withdrawal of all possible objects.

Intentionalists have answered the first two challenges. \citet{Tye1995} and \citet{Crane1998} argue that pains represent bodily disturbances at locations, and \citet{Crane1998} and \citet{Mendelovici2013} that moods are directed at the world or at everything indefinitely. Profound boredom has received less attention in this debate. We argue that it is intentional, and that its object is time itself.

\section{Framework}

\begin{definition}[Intentionality]\label{d:int}
A conscious state is \emph{intentional} if it presents something as being some way. What it presents is its \emph{object}; the way it presents it is its \emph{aspect}.
\end{definition}

The definition is Husserl's, not Brentano's: the object need not exist, need not be a particular, and need not be given as an object of attention. A mood can present the world as bleak without the world being the focus of any act.

\begin{principle}[Correctness]\label{p:corr}
If a conscious state can be correct or incorrect, apt or inapt, in virtue of how things are, it presents something as being some way.
\end{principle}

\begin{principle}[Phenomenal articulation]\label{p:art}
If a conscious state has a phenomenal character with an articulated structure, so that it is felt as a quality of something, the bearer of that quality is presented by the state.
\end{principle}

Principle~\ref{p:corr} links intentionality to normative assessment: a state can be assessed only against what it is about. Principle~\ref{p:art} links it to phenomenal character. It does not say that all phenomenal character is intentional, which would beg the question; it says that where a quality is felt as a quality of something, the something is presented. A pain felt in the foot presents the foot; a heaviness felt as the heaviness of the day presents the day.

\section{The ordinary candidates}

\begin{lemma}\label{l:sens}
Bodily sensations and undirected moods are intentional.
\end{lemma}

\begin{proof}
A pain is felt in a place, as a quality of a part of the body, and by Principle~\ref{p:art} it presents that part as being some way, namely disturbed. An undirected mood, such as elation or gloom, is felt as a quality of everything that is encountered: things look inviting or flat, the future open or closed. By Principle~\ref{p:art} it presents the world or one's situation as a whole as being some way. Both can be inapt: phantom pain presents a limb that is not there, and gloom can present as hopeless a situation that is not. By Principle~\ref{p:corr}, they are intentional.
\end{proof}

The lemma is the familiar intentionalist answer, and we do not defend it further here; PHE-21 argues that moods are pre-reflective disclosures of one's situation, fallible but not arbitrary, which entails that they are intentional in the sense of Definition~\ref{d:int}. The case that resists this answer is a state that is directed at no being, not even at the world of beings as a whole, because it is defined by their withdrawal.

\section{Profound boredom}

Heidegger distinguishes three forms of boredom \citep[\S\S 18--38]{Heidegger1995}. In the first, one is \emph{bored by} something: waiting four hours at a provincial railway station, one looks at the timetable, counts the trees, walks up and down, and keeps looking at one's watch. The station refuses to occupy us, and time drags. Heidegger identifies two structural moments: being \emph{held in limbo} (Hingehaltenheit), by the time that drags, and being \emph{left empty} (Leergelassenheit), by the things that refuse themselves. In the second form, one is \emph{bored with} something: after an evening at a pleasant dinner party, one finds that one was bored, though nothing in the evening was boring. One passed the time, and the passing of time was itself the emptiness. In the third form, profound boredom, \enquote{it is boring for one} (es ist einem langweilig), for instance walking through the streets of a city on a Sunday afternoon. Nothing in particular bores, and nothing can be done to pass the time, because every activity is equally indifferent. Beings as a whole refuse themselves.

Profound boredom thus has, on Heidegger's account, the same two structural moments as the first form, transposed. One is left empty, not by this station but by beings as a whole; one is held in limbo, not by the slowness of this afternoon but by time as such, which neither passes nor stands still. Heidegger describes the second moment as an entrancement (Gebanntsein) by the horizon of time, and says that in profound boredom Dasein is held by time itself \citep{Heidegger1995}.

\begin{lemma}\label{l:corr}
Profound boredom presents something as being some way.
\end{lemma}

\begin{proof}
Profound boredom has an articulated phenomenal character: it is felt as an emptiness and as a drag, a while that is long, as the German Langeweile says. It is felt as a quality of something, not as a free-floating sensation, since the subject can say that \enquote{it} is boring, everything and nothing. By Principle~\ref{p:art}, its bearer is presented. It is also open to assessment. A person who finds every possibility indifferent because she is exhausted may be wrong about her situation, which is in fact rich in possibilities, and come to see this when she has rested; a person in Heidegger's condition may be right that none of the activities on offer would make a difference to how she stands. By Principle~\ref{p:corr}, profound boredom presents something as being some way.
\end{proof}

\begin{lemma}\label{l:notbeing}
What profound boredom presents is not a being or an activity.
\end{lemma}

\begin{proof}
In profound boredom every being and every activity is presented as equally indifferent; that is the difference between the third form and the first two. If the state presented a particular being or activity as dragging, other beings or activities could relieve it, as a book relieves waiting at a station. Nothing relieves profound boredom, since it is indifferent to which being or activity is offered. Hence its drag is not presented as a quality of any being or activity.
\end{proof}

\begin{lemma}\label{l:time}
What profound boredom presents as dragging is time: the time of the subject's existence as such.
\end{lemma}

\begin{proof}
The drag of boredom is a temporal quality: it is a felt length of a while, the felt excess of duration over what fills it. By Lemma~\ref{l:corr}, profound boredom presents something as having this quality. By Lemma~\ref{l:notbeing}, that something is not a being or an activity, and so not the duration of a particular occupation. What remains to bear a temporal quality is time itself, not as measured by the clock, which is a being among beings, but as the time in which the subject exists: the span of her having-been and her coming, which is felt as stretched out and held still. The phenomenology confirms this. In profound boredom one does not experience a particular period as long; one experiences the present as extended without limit, not passing and not yet past, with nothing coming.
\end{proof}

\begin{theorem}\label{t:main}
Profound boredom is intentional towards time itself, not towards an activity. Hence every conscious state, including the apparent exceptions, is consciousness of something.
\end{theorem}

\begin{proof}
By Lemmas~\ref{l:corr} and~\ref{l:time}, profound boredom presents time itself as dragging and empty; by Definition~\ref{d:int}, it is intentional towards time itself, and by Lemma~\ref{l:notbeing}, not towards an activity. The other candidates are intentional by Lemma~\ref{l:sens}.
\end{proof}

\begin{remark}\label{r:ordinary}
Theorem~\ref{t:main} separates profound boredom from the boredom of tasks. Ordinary boredom, the boredom of the station or the lecture, is directed at an activity or situation and presents it as failing to engage one's capacities. EXI-10 argues that this ordinary boredom signals a mismatch between one's capacities and one's activities and has an adaptive function. Profound boredom is not a rarer case of the same state, since it has a different object; it is the limit case in which the felt drag, which in ordinary boredom is the drag of an activity, is detached from every activity and attaches to time itself.
\end{remark}

\begin{remark}\label{r:anxiety}
Profound boredom is also distinct from anxiety, though both involve the withdrawal of beings as a whole. PHE-16 argues that anxiety discloses the nothing as that withdrawal. In boredom the withdrawal is presented as indifference, and what becomes salient is the time in which one is held; in anxiety it is presented as a slipping away that leaves one without hold, and what becomes salient is the nothing. The difference is one of object, as the theorem would predict.
\end{remark}

\section{Objections and replies}

\begin{objection}[Searle's undirected anxiety]
\citet{Searle1983} distinguishes being anxious about something from undirected anxiety, a state of nervousness with no object. Such states show that some consciousness has no intentional object.
\end{objection}

\begin{reply}
Undirected anxiety is felt as a quality of one's situation as a whole: everything feels threatening or unsafe, without any particular thing being the threat. By Principle~\ref{p:art} it presents the situation as a whole. Searle's distinction is between states directed at particular objects and states that are not, and it is correct; but Definition~\ref{d:int} does not require a particular object. PHE-16 develops the case of Heideggerian anxiety, which discloses the nothing without making it an entity.
\end{reply}

\begin{objection}[Time is not an object]
Time itself is not something that can be presented. To say that boredom presents time as dragging is a metaphor for the fact that the bored subject's experience has a certain character. The adverbial description, that she experiences boredly, is the literal one, and it involves no object.
\end{objection}

\begin{reply}
Time need not be an object in the sense of an entity to be what a state presents; Definition~\ref{d:int} requires only that something be presented as being some way. The time of one's existence is presented in every experience of duration, as the living present presents its earlier phases as elapsed (PHE-05). The adverbial description cannot account for the correctness conditions of Lemma~\ref{l:corr}: a way of experiencing is not correct or incorrect unless it presents something as being some way, and the bored subject's experience can be corrected by learning that her time is not as empty as it seems.
\end{reply}

\begin{objection}[Dependence on Heidegger]
The argument relies on Heidegger's description of profound boredom. Empirical research on boredom finds a state defined by disengagement from a particular activity \citep{EastwoodEtAl2012}. Perhaps profound boredom, as Heidegger describes it, does not exist.
\end{objection}

\begin{reply}
If it does not exist, it is no counterexample to Brentano's thesis, and the theorem holds a fortiori. The argument is conditional on the description in the sense that it shows that the state, if it is as described, is intentional. The empirical research concerns the boredom of tasks, which Remark~\ref{r:ordinary} sets apart. \citet{Svendsen2005} offers a cultural history in which a diffuse, objectless boredom is a recognizable modern experience, which suggests that Heidegger's description has an object of study.
\end{reply}

\section{What the result does not establish}

The result does not establish that intentionality is the mark of the mental in the sense that it suffices for mentality: many non-mental things, such as maps, have correctness conditions. It does not settle whether the intentionality of moods and sensations is representational in the sense of the intentionalist theories of phenomenal character. And it does not establish Heidegger's further claim that profound boredom reveals the moment of decision, in which Dasein is called back to its own possibilities. The claim is only that consciousness is always consciousness of something, and that the case that seemed to show otherwise is directed at time.

\section{Conclusion}

Every conscious state presents something as being some way. Sensations present the body and moods the situation as a whole. Profound boredom, in which all beings and activities withdraw into indifference, presents what remains when they have withdrawn: the time of one's existence, as dragging and empty. It is intentional towards time itself, not towards an activity, and Brentano's thesis holds without exception.

\begin{thebibliography}{99}\small
\bibitem[Brentano(1874)]{Brentano1874} Brentano, F. (1874). \emph{Psychologie vom empirischen Standpunkte}. Duncker \& Humblot.
\bibitem[Crane(1998)]{Crane1998} Crane, T. (1998). Intentionality as the mark of the mental. In A. O'Hear (ed.), \emph{Current Issues in Philosophy of Mind}. Cambridge University Press.
\bibitem[Eastwood et~al.(2012)]{EastwoodEtAl2012} Eastwood, J.~D., Frischen, A., Fenske, M.~J. and Smilek, D. (2012). The unengaged mind: Defining boredom in terms of attention. \emph{Perspectives on Psychological Science}, 7(5), 482--495.
\bibitem[Goldie(2000)]{Goldie2000} Goldie, P. (2000). \emph{The Emotions: A Philosophical Exploration}. Oxford University Press.
\bibitem[Heidegger(1995)]{Heidegger1995} Heidegger, M. (1995). \emph{The Fundamental Concepts of Metaphysics: World, Finitude, Solitude}. Trans. W. McNeill and N. Walker. Indiana University Press.
\bibitem[Husserl(1900--1901)]{Husserl1901} Husserl, E. (1900--1901). \emph{Logische Untersuchungen}. Niemeyer.
\bibitem[Mendelovici(2013)]{Mendelovici2013} Mendelovici, A. (2013). Intentionalism about moods. \emph{Thought: A Journal of Philosophy}, 2(1), 126--136.
\bibitem[Searle(1983)]{Searle1983} Searle, J.~R. (1983). \emph{Intentionality: An Essay in the Philosophy of Mind}. Cambridge University Press.
\bibitem[Svendsen(2005)]{Svendsen2005} Svendsen, L. (2005). \emph{A Philosophy of Boredom}. Trans. J. Irons. Reaktion Books.
\bibitem[Tye(1995)]{Tye1995} Tye, M. (1995). \emph{Ten Problems of Consciousness}. MIT Press.
\end{thebibliography}

\end{document}
